% BHU PYQ -- Manually transcribed & error-corrected
% Paper: BSC-05A: Elements of Earth Science (Ancillary)
% Exam: B.Sc. (Hons.) Semester V Examination, 2013-14

\documentclass[11pt,a4paper]{article}
\usepackage[a4paper,top=2.5cm,bottom=2.5cm,left=2.8cm,right=2.8cm,headheight=14pt]{geometry}
\usepackage{amsmath,amssymb}
\usepackage[T1]{fontenc}
\usepackage[utf8]{inputenc}
\usepackage{lmodern,microtype,fancyhdr,enumitem,hyperref,lastpage,parskip}

\hypersetup{colorlinks=true,urlcolor=black,linkcolor=black,
  pdftitle={BSC-05A Elements of Earth Science (Ancillary) -- BHU B.Sc. Sem V 2013-14}}
\pagestyle{fancy}\fancyhf{}
\renewcommand{\headrulewidth}{0.4pt}\renewcommand{\footrulewidth}{0.4pt}
\fancyhead[L]{\small   \textit{Banaras Hindu University}}
\fancyhead[R]{\small   \textit{BSC-05A: Elements of Earth Science (Ancillary)}}
\fancyfoot[C]{\small Page  \thepage\ of \pageref{LastPage}}


\newlist{parts}{enumerate}{2}
\setlist[parts,1]{label=(\alph*),leftmargin=*,itemsep=5pt,topsep=3pt}

\newcommand{\pts}[1]{\hfill   \textit{[#1]}}


\begin{document}

\begin{center}
  {\large\bfseries BANARAS HINDU UNIVERSITY}\\[2pt]
  {
\normalsize B.Sc. (Hons.) --- Semester V Examination, 2013-14}\\[6pt]
  \rule{0.8\linewidth}{0.5pt}\\[6pt]
  {\large\bfseries Paper No. BSC-05A: Elements of Earth Science (Ancillary)}\\[4pt]
  {
\normalsize Time: 3 Hours \hfill Maximum Marks: 70}\\[2pt]
  {\small Ref. No.: BS/Sem V/260}
\end{center}
\rule{\linewidth}{0.4pt}
\smallskip



\noindent   \textit{Instructions: Attempt any   \textbf{five} questions, selecting at least   \textbf{two} questions each from Section A and Section B. Section C is compulsory. All questions carry equal marks (14 marks each).}
\bigskip

\bigskip


\noindent {\large\bfseries SECTION --- A}
\rule{\linewidth}{0.4pt}\smallskip




\noindent   \textbf{Question 1.} Describe the geological work of a glacier. \pts{14}

\medskip


\noindent   \textbf{Question 2.} Write a detailed note on   \textbf{any two} of the following: \pts{$2  \times 7 = 14$}

\begin{parts}
  \item Mountain building
  \item Plate Tectonics
  \item Rock cycle
\end{parts}

\medskip


\noindent   \textbf{Question 3.} Discuss the significance of Earth System science with a focus on mitigating natural disasters. \pts{14}

\medskip


\noindent   \textbf{Question 4.} What are igneous rocks? Write an essay on their classification giving suitable examples. \pts{14}

\bigskip


\noindent {\large\bfseries SECTION --- B}
\rule{\linewidth}{0.4pt}\smallskip




\noindent   \textbf{Question 5.} Discuss the various causes for the generation of earthquakes giving suitable examples. \pts{14}

\medskip


\noindent   \textbf{Question 6.} Write short notes on   \textbf{any four} of the following: \pts{$4  \times 3.5 = 14$}

\begin{parts}
  \item Cirque
  \item Attrition
  \item Chemical weathering
  \item Alluvial fan
  \item Delta
  \item Present is Key to the past
\end{parts}

\medskip


\noindent   \textbf{Question 7.} How do you distinguish between   \textbf{any two} of the following mineral pairs based on their diagnostic physical properties? \pts{$2  \times 7 = 14$}

\begin{parts}
  \item Muscovite and biotite
  \item Calcite and quartz
  \item Olivine and augite
  \item Tourmaline and hornblende
\end{parts}

\bigskip


\noindent {\large\bfseries SECTION --- C (Compulsory)}
\rule{\linewidth}{0.4pt}\smallskip




\noindent   \textbf{Question 8.} Explain the following terms in two or three sentences: \pts{$14  \times 1 = 14$}

\begin{parts}
  \item Strike
  \item Types of plate boundaries
  \item Panthalassa
  \item Schist
  \item Sandstone
  \item Paleontology
  \item Lustre
  \item Hardness in minerals
  \item Asthenosphere
  \item Syncline
  \item Net slip
  \item Exfoliation
  \item Slate
  \item Isostasy
\end{parts}

\vfill\rule{\linewidth}{0.4pt}

\begin{center}\small
\href{https://bhu-exam.vercel.app/}{Click here to visit our website}\\[4pt]
\href{https://me-aryan.vercel.app/}{Click here to visit our contributor}\\[4pt]
\href{https://quashan-me.vercel.app/}{Click here to visit our contributor}
\end{center}
\end{document}
